%!TEX root = ../thesis.tex
% Chapter: Introduction

\chapter{Introduction}

Breast cancer is the most commonly diagnosed cancer in women worldwide, with around 2.3 million new cases each year \cite{sung2021global}, and X-ray mammography is the primary tool for detecting it early, while it is still small and most treatable; large-scale reviews estimate that regular screening cuts breast-cancer mortality by around 20\% \cite{marmot2012benefits}. Reading a mammogram is time-consuming and subject to inter-reader variability, which has motivated computer-aided diagnosis (CAD) research to support the radiologist \cite{burt2018deep}.

Most modern CAD systems are built on convolutional neural networks (CNNs), which perform strongly across medical-imaging tasks \cite{burt2018deep} but face two obstacles in this setting. CNNs are data-hungry, yet public mammography datasets are small, so networks trained on them readily overfit; and they extract predominantly spatial-domain features, potentially missing the multi-scale frequency-domain texture that separates one tissue type from another. The goal is therefore to pair the CNN's learned discrimination with a stable, training-free texture descriptor that does not overfit on limited data.

The wavelet scattering transform (WST), introduced by Mallat \cite{mallat2012invariant} and Bruna and Mallat \cite{bruna2013invariant}, is well suited to this role. It cascades an image through fixed wavelet filters and modulus non-linearities to build features that are locally invariant to translation and stable to small deformations, capturing rich multi-scale, multi-orientation texture in the frequency domain. It has no parameters to learn (it is a fixed mathematical operation), so it cannot overfit and can be applied directly to small datasets.

Razali et al.\ \cite{razali2023cnn} bring these ideas together in a hybrid framework that fuses CNN and WST features and report that the fusion outperforms either branch alone. They evaluate it on the INbreast dataset \cite{moreira2012inbreast} for two contrasting tasks: classifying background tissue as fatty or fibroglandular (a property of relatively uniform texture) and classifying mass lesions as benign or malignant (a localised, structured object whose margin and internal pattern carry the diagnostic signal). This project reproduces and critically re-assesses their framework: because the paper's strongest results rest on a single small test split, a recurring theme is to re-evaluate them under a more rigorous, repeated protocol.

\section{Overview}

\textbf{Objectives.} The primary objective of this work is to reproduce, in Python, the CNN--wavelet-scattering feature-fusion results of Razali et al.~\cite{razali2023cnn} on the INbreast dataset, evaluating the three paper configurations (WS-only, CNN-only, and CNN~+~WS fusion) on both the tissue-density and benign-versus-malignant mass tasks. Beyond replication, it extends the wavelet-scattering branch along two threads: (i)~a scattering-covariance descriptor for the tissue task, and (ii)~a region-aware, spatial-map-preserving representation for the mass task. All methods are evaluated under a leakage-free, image-level protocol.

\textbf{Structure.} Chapter~\ref{sec:background} reviews the background: texture features, the wavelet scattering transform, CNNs, and feature fusion. Chapter~\ref{sec:methodology} details the dataset, preprocessing pipeline, patch extraction, the wavelet-scattering and CNN feature branches, and the classifier. Chapter~\ref{sec:results} presents the paper-style replication method by method. Chapter~\ref{sec:extensions} reports the extension experiments, led by the mass region-aware model. Chapter~\ref{sec:discussion} discusses sources of discrepancy, lessons learned, and limitations, and Chapter~\ref{sec:conclusion} concludes with directions for further work.


% ============================================================
% 2. BACKGROUND
% ============================================================
